\section{The original unsmoothed finite endpoint normalization}
\label{sec:spectral-finite-endpoint}

The scalar normalization in S25, Theorem~5.10
(\texttt{finmain}), can be tested directly.  The following result
concerns the actual Fourier Galerkin matrices of that source.
It is a fixed-window statement; it makes no assertion about
large-window selection of the lowest line or convergence of zero sets.

\begin{lemma}[AC-S40-periodic: a bounded periodic reference]
\label{AC-S40-periodic}
Translate the window to \(I=(0,\ell)\), where
\(\ell=2\log\lambda\), and put
\[
 e_n(x)=\ell^{-1/2}e^{2\pi inx/\ell},\qquad
 G_{\rm per}e_n=\tfrac12\log\bigl(1+(2\pi n/\ell)^2\bigr)e_n.
                                                        \tag{FE.1}
\]
The actual closed Weil operator satisfies
\[
                         A=G_{\rm per}+B,
                  \qquad B\in\mathcal B(L^2(I)).          \tag{FE.2}
\]
This is equality on the operator domains.  In particular periodic
translations preserve \(\operatorname{Dom}(A)\), and if \(P_N\)
selects the modes \(|n|\leq N\), then
\(P_NA(1-P_N)=P_NB(1-P_N)\) on that domain.
\end{lemma}
\begin{proof}
Let \(n(h)=e^{-|h|}/(2|h|)\), as in (RN.3).
The periodic jump kernel is
\(n(x-y)+K(x,y)\), where
\[
                     K(x,y)=\sum_{k\ne0}n(x-y+k\ell).
                                                        \tag{FE.3}
\]
For fixed \(x\in I\), tiling the complement of \(I\) gives
\(\int_I K(x,y)dy=\int_{\mathbb R\setminus I}n(x-y)dy\).
Thus the killing term in the zero-extension part form cancels
the diagonal terms in the difference of the periodic and restricted
jump forms.  It follows that
\[
                     G_I=G_{\rm per}+\mathcal K,
             \qquad (\mathcal Kf)(x)=\int_I K(x,y)f(y)dy.
                                                        \tag{FE.4}
\]
Here \(\mathcal K\) is bounded on \(L^2\).  The terms \(k=1,-1\)
are bounded by the two corner kernels
\(1/(2(x+\ell-y))\) and \(1/(2(\ell-x+y))\), respectively.
Each has norm at most \(\pi/2\), by the weighted Schur estimate
in step~3 of AC-S37.  The remaining exponentially decaying sum
is bounded on the square.

All these identities first hold with the jump singularities
truncated.  The uniform cross-kernel bounds and monotone convergence
of the nonnegative jump forms show that the two closed form domains
coincide and that their difference is exactly the bounded form
\(\langle f,\mathcal Kf\rangle\).  The associated operator identity
is therefore (FE.4), with equality of domains.
The periodic form has the diagonal multiplier (FE.1), obtained by
integrating each Fourier mode against the periodized kernel.
Finally (RN.4) says \(A=G_I+B_L\), with \(B_L\) bounded.
This proves (FE.2).  Periodic translations commute with
\(G_{\rm per}\) and preserve its domain.
\end{proof}

\begin{lemma}[AC-S40-translation: a compact-resolvent translation gain]
\label{AC-S40-translation}
Let \(\mathcal X_I\) be (RN.2), and let \(T_h\) denote periodic
translation \(T_hf(x)=f(x-h)\), with period \(\ell\).
For \(f\in\mathcal X_I\), as \(h\downarrow0\),
\[
 \frac{[A,T_h]f}{\sqrt{h\log(1/h)}}\ \rightharpoonup\ 0
                         \quad\hbox{in }L^2(I).           \tag{FE.5}
\]
The displayed vectors are bounded in norm.
Suppose that the lowest eigenvalue \(\alpha\) of \(A\) is simple,
with real normalized eigenvector \(u\).  Then
\[
             \|T_hu-u\|_2=o\bigl(\sqrt{h\log(1/h)}\bigr).
                                                        \tag{FE.6}
\]
The same conclusion holds for any \(\phi\in\mathcal X_I\)
orthogonal to \(u\) such that
\((A-\alpha)\phi=f_0\), with \(f_0\in C_c^\infty(I)\).
Consequently, for these \(u,\phi\),
\[
 \|(1-P_N)u\|_2^2+\|(1-P_N)\phi\|_2^2
                         =o(\log N/N).                   \tag{FE.7}
\]
\end{lemma}
\begin{proof}
\proofstep{1. The commutator on the full operator domain.}
Write \(E\) for zero extension, \(P\) for restriction to \(I\),
and \(S_h\) for whole-line translation.  By AC-S37 every
\(f\in\mathcal X_I\) has continuous zero extension and obeys
\[
 |f(x)|\leq C V(d(x)),\qquad
 V(d)\sim(\log(1/d))^{-1/2}
                  \quad\hbox{near the two endpoints}.    \tag{FE.8}
\]
The cross-kernel argument of step~3 of that theorem, applied at
both endpoints, gives \(Ef\in\operatorname{Dom}(G)\).
Its image equals \(G_If\) inside \(I\) and
\(-\int_I n(x-y)f(y)dy\) outside.  From (FE.8), for exterior
distance \(s\downarrow0\),
\[
                  |GEf(x)|\leq C\sqrt{\log(1/s)}.        \tag{FE.9}
\]
Indeed split the integral at distance \(s\) from the endpoint;
the remaining singular integral is bounded by
\(C\int_s^\delta dt/(t\sqrt{\log(1/t)})\), and the rest is bounded.

For \(0<h<\ell/4\), put
\[
 v_h=S_{h-\ell}Ef,\qquad w_h=1_{(0,h)}v_h,
 \qquad ET_hf=S_hEf+(1-S_\ell)w_h.                        \tag{FE.10}
\]
The same bounded cross-operator decomposition shows that restriction
to an interval followed by zero extension preserves
\(\operatorname{Dom}(G)\).  Thus all terms in (FE.10) are in that
domain.  This assertion can equivalently be checked by truncating
\(n\), splitting the form at the interval endpoints, and passing
to the limit with the uniform Carleman bound; the off-diagonal
blocks are bounded operators, so both the form and operator
decompositions are valid.

We estimate \(Gw_h\) on and off \((0,h)\).  Off this interval
it is minus the cross-kernel applied to \(w_h\), and hence
\[
 \|1_{\mathbb R\setminus(0,h)}Gw_h\|_2
              \leq C\|w_h\|_2\leq C\sqrt h\,V(h).        \tag{FE.11}
\]
On \((0,h)\), it equals
\[
 Gv_h(x)+\int_0^{\ell-h}n(x+t)f(\ell-h-t)dt.             \tag{FE.12}
\]
The first term is a translate of the bounded function \(G_If\).
For the integral, its part with \(t<h\) is bounded by
\(CV(h)\log(1+h/x)\), whose \(L^2(0,h)\) norm is
\(O(\sqrt h V(h))\).  For \(h<t<\delta\), use
\(V(h+t)\leq C V(t)\) and integrate \(CV(t)/t\).
This gives \(O(\sqrt{\log(1/h)})\) pointwise.
The remaining part is bounded.  Therefore
\(Gw_h/\sqrt{h\log(1/h)}\) is bounded in \(L^2(\mathbb R)\)
and converges weakly to zero: its part on \((0,h)\) has shrinking
support, while its complementary part tends strongly to zero
by (FE.11).

Using (FE.10) and commutation of \(G\) with \(S_h\),
\[
 [G_I,T_h]f
  =PS_hGEf-T_hG_If+PG(1-S_\ell)w_h.                      \tag{FE.13}
\]
The first difference is supported on \((0,h)\).
Its norm is \(O(\sqrt{h\log(1/h)})\) by boundedness of \(G_If\)
and (FE.9), so its normalized value tends weakly to zero as well.
This proves (FE.5) for \(G_I\).

By (RN.4), \(A-G_I\) is a scalar plus a compressed convolution
by a finite signed measure; the kernel \(k=\mu-n\) is integrable
and the prime/polar terms have finite total variation on a fixed
window.  Applying (FE.10) to that convolution shows that its
commutator has norm \(O(\sqrt h)+O(\sqrt h V(h))\): the wrapping
strip contributes the first term, and the two \(w_h\) terms the
second.  Dividing by \(\sqrt{h\log(1/h)}\) gives strong convergence
to zero.  This establishes (FE.5) for the actual \(A\).

\proofstep{2. Compactness turns weak smallness into a strict gain.}
Let \(Q=1-|u\rangle\langle u|\), and let
\(R=Q(A-\alpha)^{-1}Q\) be the reduced resolvent.
It is compact.  Equations (FE.2) and (FE.5) justify
\[
 Q(T_hu-u)=R[A,T_h]u
                  =o\bigl(\sqrt{h\log(1/h)}\bigr).        \tag{FE.14}
\]
Since \(u\) is real, \(r_h=\langle u,T_hu\rangle\) is real
and tends to one.  Unitarity gives
\(\|T_hu-u\|_2^2=2(1-r_h)\) and
\(\|QT_hu\|_2^2=1-r_h^2\); their ratio tends to one.
This proves (FE.6).

For \(\phi\) in the statement,
\((A-\alpha)(T_h\phi-\phi)=(T_h-1)f_0+[A,T_h]\phi\).
The first term is \(O(h)\) in \(L^2\), and the same compactness
argument controls its \(Q\)-component.
The remaining scalar component is bounded by
\(\|T_{-h}u-u\|_2\|\phi\|_2\); here
\(\|T_{-h}u-u\|_2=\|T_hu-u\|_2\).
Thus (FE.6) also holds for \(\phi\).

\proofstep{3. Fourier tails.}
Average \(\|T_hf-f\|_2^2\) over \(N^{-1}<h<2N^{-1}\).
For \(|n|>N\), the average of
\(|1-e^{-2\pi inh/\ell}|^2\) is bounded below by a positive
constant depending only on \(\ell\).
To see this, rescale \(h=t/N\); the resulting continuous
function of \(|n|/N\geq1\) is strictly positive and tends to two
at infinity.  The averaged little-oh bound from (FE.6)
therefore proves (FE.7).
\end{proof}

\begin{theorem}[AC-S40: decay of the actual finite endpoint and failure of its unrescaled determinant limit]
\label{AC-S40}
Let \(r_0>1\) solve \(r_0^3-r_0-1=0\), and fix
\[
                           1<\lambda\leq\sqrt{r_0}.       \tag{FE.15}
\]
Let \(u_N\) be the \(L^2\)-normalized lowest eigenvector of the
actual S25 Galerkin matrix, with phases converging to the nonnegative
normalized ground state \(u\).  The finite lowest eigenvalues are
eventually simple and even, and
\[
 0<|u_N(0)|=|u_N(\ell)|
                       \leq \frac{C_\lambda}{\sqrt{\log N}}
                       \longrightarrow0.                \tag{FE.16}
\]
Here the coordinates are logarithmic on \((0,\ell)\),
\(\ell=2\log\lambda\), so both endpoint values equal the
source functional \(\delta_N(u_N)\).

For the exact source normalization
\(\xi_N=u_N/u_N(\ell)\), the regularized determinant
in S25, Theorem~5.10 (\texttt{finmain}), consequently satisfies
\[
 \bigl|\det{}_{\rm reg}(D^{(\lambda,N)}_{\log}-0)\bigr|
       =|\widehat\xi_N(0)|\ \geq c_\lambda\sqrt{\log N}
                       \longrightarrow\infty.            \tag{FE.17}
\]
Thus these unrescaled determinants do not have a finite locally
uniform entire limit on this range.  Multiplication by a suitable
nonzero scalar is essential, as in AC-S39.  No conclusion about
divergence of their zero sets follows from (FE.17).
All statements are unconditional.
\end{theorem}
\begin{proof}
\proofstep{1. An odd test with nonzero renormalized trace.}
Write \(A=A_L\) in translated coordinates.  AC-S28 supplies a
simple even lowest line, with \(u\geq0\), and AC-S06 gives
\(u_N\to u\) in \(L^2\) and form norm, \(\alpha_N\to\alpha\),
and eventual even simplicity.  The real phase can be chosen so
that this convergence holds.

We need one fixed imaginary odd vector \(\phi=iw\), where \(w\)
is real and odd about \(\ell/2\), satisfying
\[
 (A-\alpha)\phi=f_0\in iC_c^\infty(I;\mathbb R),\quad
 f_0\hbox{ odd},\qquad \tau_+(w)=:\tau\ne0.               \tag{FE.18}
\]
Here is a construction which does not assume nonvanishing for an
arbitrary odd eigenfunction.  The compact barrier in AC-S38,
placed at the left endpoint and reflected at the right endpoint,
has bounded \(A\)-image, by the calculation (RN.20) and the
bounded perturbation (RN.4).  Their imaginary odd difference
\(\phi_0\) therefore belongs to \(\mathcal X_I\) and has nonzero
right trace.  Its forcing \((A-\alpha)\phi_0\) is bounded and odd.

The odd resolvent \(R_o=(A_{\rm odd}-\alpha)^{-1}\) maps
\(L^\infty_{\rm odd}\) boundedly into \(L^\infty_{\rm odd}\).
For precision, use its semigroup integral.  On \(0<t<1\),
the finite-measure Dyson bound in AC-S36 gives a uniform
\(L^\infty\) bound.  On \(t\geq1\), first apply
\(e^{-A}:L^2\to L^\infty\) and then use the positive gap
\(\alpha_{\rm odd}-\alpha>0\); this gives an integrable
exponential bound after multiplication by \(e^{t\alpha}\).
The integral agrees with the \(L^2\) resolvent by the spectral
theorem.

Approximate \((A-\alpha)\phi_0\) in \(L^2\) by imaginary odd
smooth compact functions \(f_j\), with a uniform \(L^\infty\)
bound, and put \(\phi_j=R_of_j\).
Then \(\phi_j\to\phi_0\) and \(G_I\phi_j\to G_I\phi_0\)
in \(L^2\), while both sequences are uniformly bounded in
\(L^\infty\).  Formula (RN.15) implies convergence of their
traces: both weights there are in \(L^1\), and an \(L^2\)-convergent,
uniformly bounded sequence converges against every \(L^1\) weight,
by truncating that weight to an \(L^2\) function.
Choose one \(j\) with nonzero trace and set \(\phi=\phi_j\).
This proves (FE.18).

\proofstep{2. Quantitative Ritz remainders.}
Apply (FE.7) to \(u\) and \(\phi\).
Abel summation of the tail function
\(\sum_{|n|>M}|u_n|^2=o(\log M/M)\) gives
\[
 \sum_{|n|>N}\log(2+|n|)|u_n|^2=o(\log^2N/N).
\]
Indeed the first boundary term is
\(\log N\,o(\log N/N)\), and the remaining integral is
bounded using \(\int_N^\infty (\log x)x^{-2}dx\).
By (FE.2), the Rayleigh quotient of \(P_Nu\) yields
\[
                 0\leq\alpha_N-\alpha=o(\log^2N/N).
                                                        \tag{FE.19}
\]
The exact identity used here is
\(q_A(P_Nu)-\alpha\|P_Nu\|_2^2
 =q_A((1-P_N)u)-\alpha\|(1-P_N)u\|_2^2\).
The Fourier cutoffs are a form core by (FE.2).

\proofstep{3. The exact displacement vector and its Abelian growth.}
S25, Lemma~5.1 (\texttt{basicexpli}), (\texttt{form}),
Lemma~5.2 (\texttt{basics}), (\texttt{QD}), and
Lemma~5.4 (\texttt{key}), give, for \(D_Ne_n=ne_n\),
\[
 \begin{gathered}
 (A_N-\alpha_N)D_Nu_N=-c_N\beta_N,\qquad
 c_N=\langle\eta_N,u_N\rangle=\sqrt\ell\,u_N(\ell),\\
 \eta_N=\sum_{|n|\leq N}e_n,\qquad
 \beta_N=\sum_{|n|\leq N}b_ne_n.
 \end{gathered}                                          \tag{FE.20}
\]
The source's argument between Definition~5.3
(\texttt{even-simple}) and Lemma~5.4 proves \(c_N\ne0\)
under eventual even simplicity.  It does not give its size.

There are no prime translations in (FE.15), because
\(e^\ell\leq r_0<2\).  The formula for \(b_n\) in the proof
of \texttt{basicexpli}, or directly the off-diagonal part of
(EP.5)--(EP.6), is
\[
 b_n=\frac1\pi\int_0^\ell
       \left(\frac{e^{-h/2}}{1-e^{-2h}}-2\cosh(h/2)\right)
                         \sin(2\pi nh/\ell)\,dh
       =\tfrac14\operatorname{sgn}(n)+O_\ell(1/|n|).
                                                        \tag{FE.21}
\]
The last equality follows by subtracting \(1/(2h)\);
the remainder extends smoothly to \(h=0\), so integration by
parts gives \(O(1/|n|)\), while
\(\int_0^\ell\sin(2\pi nh/\ell)dh/h
 =\operatorname{Si}(2\pi n)\to(\pi/2)\operatorname{sgn}(n)\).
In particular \(\|\beta_N\|_2=O(\sqrt N)\).

The coefficients of \(\phi=iw\) are real and odd, and
\(\phi_n=\ell^{-1/2}\int_0^\ell w(x)\sin(2\pi nx/\ell)dx\).
By (RN.1) and (FE.18),
\(w(x)\sim-\tau/\sqrt{\log(1/x)}\) near zero.
The exact identity
\[
 \sum_{n=1}^N\sin(n\theta)
       =\frac{\cos(\theta/2)-\cos((N+1/2)\theta)}
                    {2\sin(\theta/2)}                  \tag{FE.22}
\]
therefore gives
\[
 \sum_{n=1}^N\phi_n
       =-\frac{2\tau\sqrt\ell}{\pi}\sqrt{\log N}
                               +o(\sqrt{\log N}),
 \qquad
 \langle\beta_N,\phi\rangle
       =-\frac{\tau\sqrt\ell}{\pi}\sqrt{\log N}
                               +o(\sqrt{\log N}).        \tag{FE.23}
\]
We detail the limiting estimate to avoid a pointwise Fourier
convergence assumption.  Reflection reduces the integral to twice
the left half interval.  For a fixed small \(\delta\), its interior
part is bounded uniformly in \(N\), by (FE.22).
Near zero replace \(2\sin(\pi x/\ell)\) by \(2\pi x/\ell\)
and \(\cos(\pi x/\ell)\) by one; the errors have bounded integrals.
The leading integral is
\[
 \int_0^\delta\frac{1-\cos((N+1/2)2\pi x/\ell)}
                         {x\sqrt{\log(1/x)}}dx
                         \sim2\sqrt{\log N}.             \tag{FE.24}
\]
Below \(1/N\) it is \(O((\log N)^{-1/2})\).
Above \(1/N\), the nonoscillatory integral is
\(2\sqrt{\log N}+O(1)\); integration by parts or Dirichlet's
test bounds its cosine part by \(O((\log N)^{-1/2})\), since
\(1/(x\sqrt{\log(1/x)})\) is decreasing on this small interval.
The \(o((\log(1/x))^{-1/2})\) remainder in \(w\) contributes
\(o(\sqrt{\log N})\), using
\(C\min(N,1/x)\) as a bound for (FE.22), then first fixing
\(\delta\) and finally decreasing it.
This proves the first formula in (FE.23) with its constant.
The \(O(1/|n|)\) remainder in (FE.21) has uniformly bounded
pairing with \(\phi\) by Cauchy--Schwarz, and the two odd
Fourier halves give the second formula.

\proofstep{4. The finite endpoint tends to zero.}
The vector \(D_Nu_N\) is odd.  The odd spectral gap and
\(\alpha_N\to\alpha\), together with (FE.20), give
\[
                         \|D_Nu_N\|_2\leq C|c_N|\sqrt N.
                                                        \tag{FE.25}
\]
Pair (FE.20) against \(P_N\phi\).  By (FE.2) and (FE.18),
\[
 (A_N-\alpha_N)P_N\phi
       =P_Nf_0+(\alpha-\alpha_N)P_N\phi
                                  -P_NB(1-P_N)\phi.     \tag{FE.26}
\]
The absolute value of the pairing with the first term is
\(|\langle D f_0,u_N\rangle|\leq\|D f_0\|_2\), using an inner
product linear in its second argument.
Here \(D=(\ell/(2\pi i))\partial_x\), and \(f_0\) is smooth
and compactly supported.  The absolute value of the remaining
pairing is at most
\[
 C|c_N|\sqrt N
   \left(|\alpha_N-\alpha|\,\|\phi\|_2
                     +\|B\|\,\|(1-P_N)\phi\|_2\right)
                         =o(|c_N|\sqrt{\log N})           \tag{FE.27}
\]
by (FE.7) and (FE.19).
The nonzero leading coefficient in (FE.23) thus implies
\(|c_N|\sqrt{\log N}\leq C\) for all sufficiently large \(N\).
Together with the exact nonvanishing in (FE.20), this is (FE.16).

\proofstep{5. The exact source determinant.}
Since \(u\geq0\) and \(\|u\|_2=1\), its integral is positive.
The \(L^2\) convergence gives
\(\int_Iu_N\to\int_Iu>0\).  Hence
\[
 |\widehat\xi_N(0)|
   =\frac{|\int_Iu_N(x)dx|}{|u_N(\ell)|}
                              \geq c_\lambda\sqrt{\log N}.
\]
Theorem~5.10 (\texttt{finmain})(ii) gives exactly
\(\det_{\rm reg}(D^{(\lambda,N)}_{\log}-z)
 =-i\lambda^{-iz}\widehat\xi_N(z)\).
At zero the phase has modulus one, proving (FE.17).
This establishes the failure of the particular unrescaled scalar
limit, while leaving zero locations and the scalar-renormalized
limit of AC-S39 untouched.
\end{proof}

\paragraph{Scope of the remaining spectral question.}
On (FE.15), ordinary finite endpoint normalization is now decided:
it tends to zero on normalized vectors, and its reciprocal makes
the source determinants diverge at zero.
AC-S41 below determines the nonzero limit of
\(\sqrt{\log N}\,u_N(\ell)\), and AC-S42 extends that conclusion
to every fixed window with a simple even minimum, including all
prime displacement terms.  The arbitrary-window simple-even
property itself remains unproved.  None of these fixed-window conclusions
proves the large-window ground-line selection demanded in S25.
